\section{Causal simulation and finite execution}\label{sec:resources}
The preceding cardinalities count persistent labels. A programme describing real-valued functions or arbitrarily long tables is not free bit complexity. We separate statistical comparison, numerical execution, and one finite class where synthesis is effective.

\begin{proposition}[Typed common-task transfer]\label{prop:simulation}
Let $\mathfrak E$ and $\mathfrak F$ be prepared experiments with the same bounded scored task, $0\le L\le H$. Suppose a parameter-independent causal simulator implements every legal target strategy in a specified class, including its action interface, reports, stopping rule, and task readout. Assume it uses at most $R$ internal labels, transforms the target acquisition bound $N$ to a declared physical bound $C(N)$, and gives total variation at most $\varepsilon$ between the complete stopped and scored laws, uniformly in the compared preparations and strategies. Then
\begin{equation}\label{eq:simulation}
 \mathcal R_{\mathfrak E}(C(N),RM)
 \le \mathcal R_{\mathfrak F}(N,M)+H\varepsilon.
\end{equation}
Two such simulators compose with the product of their internal cardinalities, composition of their acquisition/physical-time maps, and the sum of their complete-law TV defects. A reverse risk lower requires a reverse certificate covering every strategy in the class to which the lower is applied.
\end{proposition}
\begin{proof}
For each target observer store the pair of its current label and the simulator label. The causal strategy lift supplies only legal actions, transforms reports at the promised completion times, and uses the same scored readout; the direct persistent cardinality is at most $RM$. A bounded loss changes by at most $H$ times total variation, so taking infima proves \eqref{eq:simulation}. For composition, insert the intermediate simulated law, use contraction of total variation under the next parameter-independent kernel, and then its own defect bound. The product-state and composed-call accounts follow from the direct implementation. Without a reverse strategy lift this construction says nothing about source strategies not obtained from a target observer, which proves the stated converse qualification.
\end{proof}
Programme descriptions concatenate, scratch workspaces add in a direct simultaneous implementation (or can be reused when their lifetimes are proved disjoint), and arithmetic and physical execution times add along the actual calls. An externally supplied stage index is not automatically supplied by a simulator of an internally timed task. The proposition compares total risks of a common task; subtracting two different full-history baselines would not preserve the inequality. Marginal report TV, an unscored transcript comparison, or a fixed-time simulator without stopping compatibility is insufficient. No large-deviation equivalence is inferred from this proposition.

\begin{proposition}[Certified execution error]\label{prop:numerical}
Consider an ideal finite observer for a task bounded by $H$ and at most $n$ physical transitions. Suppose an implementation has a coupling with the ideal process such that, while the complete pasts agree, the next physical/label transition fails to agree with conditional probability at most $e_t$. A bound under the ideal prefix law on the corresponding first-mismatch event is also sufficient. If the final readout perturbation changes the conditional task loss by at most $\omega(v)$, then
\begin{equation}\label{eq:numerical}
 R_{\mathrm{impl}}\le R_{\mathrm{ideal}}
    +H\min\{1,\sum_{t=0}^{n-1}e_t\}+\omega(v).
\end{equation}
An ideal observer using approximate Bellman selections adds its actual expected action slacks $\sum_t A_t$ as in \eqref{eq:slack}; these slacks need not be counted a second time in $e_t$.
\end{proposition}
\begin{proof}
Partition the event that the paths disagree by the first transition at which they do. Before that time both constructions have the ideal prefix, so the probability is at most $\sum_te_t$, capped at one. On the complementary event the scored outcome is the same apart from the asserted readout change. The loss bound proves the inequality. The action-slack statement is the exact telescope, not a stability estimate between separately optimized controllers.
\end{proof}
For the grids of Section~\ref{sec:geometry}, a coordinate error at most $b_t$ can change a cell only in slabs of total ideal acquired mass at most
$\min\{1,C_d\rho b_t/h_t\}$, for $b_t\le h_t$, plus a separately certified exceptional evaluation probability. Coordinates mean the fixed affine simplex or traceless-matrix chart used for volume and grid cells. These are unconditional ideal masses; conditioning on survival of earlier numerical comparisons would bias them. Certified kernel couplings, integration/barycenter routines, Bellman selection errors, report precision, programme bits, scratch space and time must be supplied to apply this observation. It is not an algorithm for arbitrary Borel inputs. Predictive-coordinate error used to choose a cell is distinct from a Brier probability-readout error, whose loss modulus on a bounded probability simplex is linear in its Euclidean size.

\begin{theorem}[Effective finite rational compilation]\label{thm:finite_compiler}
Suppose the physical state, action, report and terminal-decision sets are finite, with cardinalities $S,A,Y,U$, respectively. Preparation, phase-indexed raw transition probabilities and bounded losses are rational, and $n,M$ are fixed integers. For the stationary machine convention of Definition~\ref{def:machines}, including the phase-dependent physical legality constraints, nonemptiness and the optimum risk are decidable; whenever nonempty the optimum is attained by a machine with real algebraic row probabilities. No efficiency bound for the synthesis is asserted.

For every integer $Q\ge1$, this machine admits a rational, support-preserving rounding with denominator $Q$, at the same $M$, which is still legal and stops exactly at $n$, with risk increase at most
\begin{equation}\label{eq:compiler_error}
 H\min\{1,n(A+M)/Q\}.
\end{equation}
It can be stored using
\[
 O\big((MA+M^2AY)\log(Q+1)+M\log(U+1)+M\big)
\]
programme bits, in addition to the given model and a fixed interpreter. Its transition scratch space is
$O(\log(Q+1)+\log(M+1)+\log(A+1)+\log(Y+1))$ bits. Exact rational sampling has finite expected runtime; limiting each sample to $k$ rejection trials yields a bounded runtime and an additional risk error at most $2Hn2^{-k}$, without violating legality or the deadline. The program-size statement is a feasible upper, not an optimal description-length law.
\end{theorem}
\begin{proof}
Enumerate the finite choices of initial label, stopping subset and terminal decisions. For one such choice the unknowns are the stationary probabilities $\alpha_z(a)$ and $k(j\mid z,a,y)$, each in a finite simplex. Enumerating length-at-most-$n$ physical/label paths expresses their probabilities as polynomials with rational coefficients. Requiring zero probability of every premature stop, illegal action and non-stop at $n$ gives polynomial equalities. Each forbidden-path probability is a sum of nonnegative monomials. Expected loss is a polynomial on this closed subset of the compact product of simplexes. Thus a minimum exists when the set is nonempty. Quantifier elimination over real closed fields decides feasibility, gives its optimal algebraic value and an algebraic minimizing point; take the minimum over the finitely many enumerated choices\cite{bpr}. This procedure may be extremely expensive.

For each probability row of length $r$, round to integer numerators totaling $Q$ without introducing mass outside its original support. Taking floors and distributing the remaining units only within that support gives total variation at most $r/Q$. Coupling the action and label rows at each step gives mismatch probability at most $(A+M)/Q$. Crucially, support-preserving rounding cannot create any forbidden physical/label path: every such path has a zero raw factor or a zero original controller factor, and these zero factors remain zero. Hence legality and exact stopping hold with probability one for the rounded machine, not merely approximately. Proposition~\ref{prop:numerical} proves \eqref{eq:compiler_error}.

Store one action row per label and one next-label row for each $(z,a,y)$, together with stopping flags and decisions. This gives the displayed bit bound. To sample a denominator-$Q$ row, draw $\lceil\log_2Q\rceil$ fresh bits and reject values outside $\{0,\ldots,Q-1\}$. Acceptance probability is greater than one half, and cumulative sums find the selected outcome by scanning the row. Only an index, a sum and the sampled integer are needed in scratch space. The expected number of trials is less than two; row search uses $O((A+M)\log(Q+1))$ elementary bit operations per paired action/transition sample, in addition to the raw physical call. After $k$ rejections use a fixed member of the same support. Each fallback has probability at most $2^{-k}$, so two samples per call add at most $2n2^{-k}$ to the coupling defect. The bounded-trial variant additionally uses $O(\log(k+1))$ scratch bits for an erased trial counter. The number of sampling bit operations is then bounded by $O(n(k+A+M)\log(Q+1))$. The fallback never introduces a new support edge and therefore preserves the exact interface. A trial counter is scratch space erased before the next physical call, not a persistent random seed.
\end{proof}
The general continuous-instrument results do not inherit this finite rational decidability. Conversely, the finite theorem does not replace the main compact-carrier criterion by an enumeration of one special model.

\begin{corollary}[Attained acquisition, memory and error coordinates]\label{cor:four_resources}
Take a regular engine of Theorem~\ref{thm:regular}, or a blind engine of Theorem~\ref{thm:blind}. Before its $n\ge1$ calls, make one paid early call observing an independent fair bit through an erasure channel with erasure probability $2\varepsilon$, $0\le\varepsilon\le1/2$. An estimate of this bit must be written immediately to a write-only scoring port. Also prepare an independent offset taking $\pm\delta$ with equal probability, $0\le\delta\le1$, never reported to the predictor. Its squared prediction error is scored together with the engine's loss and the early bit error. One final paid audit is performed by the scoring authority after the predictor has halted and fixed all decisions, without releasing information beforehand. Thus $N=n+2$ counts physical acquisitions and the predictor halts after $N-1$ calls; it does not make a further transition for the audit.

The source-to-perfect-bit causal deficiency at the early completion deadline is exactly $\varepsilon$, uniformly over the two bit values. For the autonomous class, all $M\ge N$ satisfy, in the regular case,
\begin{equation}\label{eq:four_regular}
 \mathcal R(N,M,\delta,\varepsilon)-B_{N-2}^*
 \asymp (M-N+1)^{-2/d}+\delta^2+\varepsilon,
\end{equation}
and, in the blind case, the exact identity is
\begin{equation}\label{eq:four_blind}
 \mathcal R(N,M,\delta,\varepsilon)
 =q_{\lfloor(M-2)/(N-2)\rfloor}(\mu)+\delta^2+\varepsilon/2.
\end{equation}
The constants in \eqref{eq:four_regular} are those of the fixed regular class, independent of these budgets and errors. The calibration and deficiency terms are actual inaccessible alternatives and actual interface defects, not arbitrary permitted error tolerances.
\end{corollary}
\begin{proof}
The best bit prediction on erasure is $1/2$, giving squared loss $\varepsilon/2$; the offset has conditional mean zero and variance $\delta^2$. These variables are independent of the engine and scored additively under one product law, so their lower bounds add to its risk for any observer. The physical interface permits only the early action at the initial label and only engine actions afterwards. That initial label is therefore distinct from all engine labels. The write-only output is not readable and cannot function as a persistent register. Any extra bit-dependent randomization of the engine can be generated as fresh independent initial randomization; it supplies no information about the independent engine variable and cannot improve the minimum over initial engine programs. Conditional on that seed, the corresponding legal engine begins in one of at most $M-1$ labels. Taking the best such conditional program proves the engine lower with $M-1$. The matching upper uses one early label and the corresponding engine machine with $M-1$ labels; the bit estimate is emitted during the early transition. After the predictor halts and fixes its decisions, the scoring authority performs the final physical audit; this adds one charged acquisition but no predictor transition. Thus $N$ counts physical acquisitions, while the autonomous predictor itself halts after $N-1$ calls. This distinction is part of the stated task, not a free terminal signal during the engine. Applying the two engine laws gives the displayed formulas.

For the deficiency, the two erased-bit source report laws have TV distance $1-2\varepsilon$, whereas the perfect-bit laws have distance one. Any common parameter-independent postprocessing preserves the former upper bound. If its error from each perfect law were at most $e$, the triangle inequality would give $1\le2e+1-2\varepsilon$, hence $e\ge\varepsilon$. Reporting the observed bit and guessing fairly on erasure attains error $\varepsilon$ for both values. The imposed early deadline forbids using the later audit. Physical preparation, erasures, engine calls and the final scoring call are all counted; the calibration offset is not an unknown kernel that may be learned from hidden free samples.
\end{proof}
